\documentclass[11pt,a4paper]{article}
% Shared preamble for RuPhO English problem statements (match T1 2026 style).
% Use: \input{latex/rupho_en_preamble.tex} after \documentclass.
\usepackage[margin=1in]{geometry}
\usepackage{amsmath,amssymb}
\usepackage{graphicx}
\usepackage{float}
\usepackage{enumitem}
\usepackage{microtype}
\usepackage{caption}
\usepackage{tabularx}
\usepackage{hyperref}

\captionsetup{font=small,labelfont=bf,skip=6pt}

\setlist[itemize]{leftmargin=1.2cm}
\setlist[enumerate]{leftmargin=1.2cm}

% One outer box per subproblem: [id | statement | points] (no inner rules)
\newcommand{\problembox}[3]{%
  \par\addvspace{0.42em}%
  \noindent
  \setlength{\fboxsep}{7.5pt}%
  \setlength{\fboxrule}{0.95pt}%
  \fbox{%
    \begin{minipage}{\dimexpr\linewidth-2\fboxsep-2\fboxrule\relax}
    \begin{tabularx}{\linewidth}{@{}>{\raggedright\arraybackslash\bfseries}p{2.1em} @{\hspace{0.9em}} >{\raggedright\arraybackslash}X @{\hspace{0.9em}} >{\raggedleft\arraybackslash\bfseries}p{2.4em}@{}}
    #1 & #2 & #3
    \end{tabularx}
    \end{minipage}%
  }%
  \par\addvspace{0.32em}%
}


\title{\textbf{Thermodynamics of ferroelectrics}}
\author{\normalsize Y24 (2024) --- English translation}
\date{}
\begin{document}
\maketitle

Ferroelectrics are substances that, at a sufficiently low temperature, have non-zero electric polarization in a zero external electric field. In this respect, ferroelectrics are similar to ferromagnets, which have a non-zero magnetic moment without an applied magnetic field. This problem will consider the thermodynamic properties of ferroelectrics, in particular the change in temperature with a change in the electric field (the so-called electrocaloric effect).  This effect allows the use of ferroelectric films for surface cooling (see figure).

Ferroelectrics are substances that, at a sufficiently low temperature, have non-zero electric polarization in a zero external electric field. In this respect, ferroelectrics are similar to ferromagnets, which have a non-zero magnetic moment without an applied magnetic field. This problem will consider the thermodynamic properties of ferroelectrics, in particular the change in temperature with a change in the electric field (the so-called electrocaloric effect).  This effect allows the use of ferroelectric films for surface cooling (see figure).

\begin{figure}[H]
    \centering
    \includegraphics[width=0.72\textwidth]{T4_EN_assets/fig1.jpg}
    \caption{Figure 1.}
\end{figure}

The state of a ferroelectric is characterized by polarization $\vec{P}$ (dipole moment per unit volume of its substance), as well as by the electric induction vector $\vec{D}$ and electric field strength $\vec{E}$. Further, in the entire problem, we consider a flat ferroelectric sample, the polarization of which is uniform and always directed along the same axis, therefore, instead of vectors, we can use the scalar quantities $P$, $D$, $E$ – projections onto this axis. Consider it known that the change in the volumetric density of electrical energy (including the energy of the electric field and the energy of the polarized matter) has the form $$ dw = E d D. $$Будем расcmатривать сегнетоэлектрик постоянного объема $V$, так что он не совершает механической работы. Тогда изменение его внутренней энергии задается выражением $$ dU = \textbackslash\{\}delta Q + V E dD, $$где $\delta Q$ – подведенное в процессе тепло.

Математическое дополнение Расcmотрим функцию двух переменных $f(x,y)$.  Частной производной функции называется производная этой функции по одному из ее аргументов при фиксированном втором аргументе. При вычислениях в термодинамике важно помнить не только по какой переменной производится дифференцирование, но и какая величина остается постоянной. Эту величину обозначают нижним индексом. Так частная производная функции $f$ по $x$ при фиксированном $y$ обозначается как $$\textbackslash\{\}left( \textbackslash\{\}frac\{\textbackslash\{\}partial f\}\{\textbackslash\{\}partial x\}\textbackslash\{\}right)_y.$$ Значение производной зависит от того, какая именно величина остается постоянной.  В качестве примера расcmотрим некоторое количество идеального  одноатомного газа, давление которого является функцией объема и давления. Тогда производные давления по объему при постоянной температуре и при постоянной энтропии $S$ (т.е. в адиабатическом процессе) будут различаться: $$\textbackslash\{\}left(\textbackslash\{\}frac\{ \textbackslash\{\}partial p\}\{\textbackslash\{\}partial V\}\textbackslash\{\}right)_T \textbackslash\{\}neq \textbackslash\{\}left(\textbackslash\{\}frac\{ \textbackslash\{\}partial p\}\{\textbackslash\{\}partial V\}\textbackslash\{\}right)_S .$$ Пусть некоторая величина $z = z(x,y)$ является функцией двух других. Из этого уравнения можно выразить величины $x = x(y,z)$ и $y = y(x,z)$. Будем считать, что такое выражение единственно. Тогда существует следующее соотношение между частными производными: \textbackslash\{\}[\textbackslash\{\}left( \textbackslash\{\}frac\{\textbackslash\{\}partial x\}\{\textbackslash\{\}partial y\}\textbackslash\{\}right)_z \textbackslash\{\}left( \textbackslash\{\}frac\{\textbackslash\{\}partial y\}\{\textbackslash\{\}partial z\}\textbackslash\{\}right)_x \textbackslash\{\}left( \textbackslash\{\}frac\{\textbackslash\{\}partial z\}\{\textbackslash\{\}partial x\}\textbackslash\{\}right)_y = -1. \textbackslash\{\}tag\{1\} \textbackslash\{\}] Отметим также, что при дифференцировании при тех же зафиксированных переменных справедливо равенство $$\textbackslash\{\}left( \textbackslash\{\}frac\{\textbackslash\{\}partial x\}\{\textbackslash\{\}partial y\}\textbackslash\{\}right)_z = \textbackslash\{\}frac\{1\}\{\textbackslash\{\}left( \{\textbackslash\{\}partial y\}/\{\textbackslash\{\}partial x\}\textbackslash\{\}right)_z\}.$$

Расcmотрим функцию двух переменных $f(x,y)$.  Частной производной функции называется производная этой функции по одному из ее аргументов при фиксированном втором аргументе. При вычислениях в термодинамике важно помнить не только по какой переменной производится дифференцирование, но и какая величина остается постоянной. Эту величину обозначают нижним индексом. Так частная производная функции $f$ по $ x$ при фиксированном $y$ обозначается как $$\textbackslash\{\}left( \textbackslash\{\}frac\{\textbackslash\{\}partial f\}\{\textbackslash\{\}partial x\}\textbackslash\{\}right)_y.$$ Значение производной зависит от того, какая именно величина остается постоянной.  В качестве примера расcmотрим некоторое количество идеального  одноатомного газа, давление которого является функцией объема и давления. Тогда производные давления по объему при постоянной температуре и при постоянной энтропии $S$ (т.е. в адиабатическом процессе) будут различаться: $$\textbackslash\{\}left(\textbackslash\{\}frac\{ \textbackslash\{\}partial p\}\{\textbackslash\{\}partial V\}\textbackslash\{\}right)_T \textbackslash\{\}neq \textbackslash\{\}left(\textbackslash\{\}frac\{ \textbackslash\{\}partial p\}\{\textbackslash\{\}partial V\}\textbackslash\{\}right)_S .$$

Пусть некоторая величина $z = z(x,y)$ является функцией двух других. Из этого уравнения можно выразить величины $x = x(y,z)$ и $y = y(x,z)$. Будем считать, что такое выражение единственно. Тогда существует следующее соотношение между частными производными: \textbackslash\{\}[\textbackslash\{\}left( \textbackslash\{\}frac\{\textbackslash\{\}partial x\}\{\textbackslash\{\}partial y\}\textbackslash\{\}right)_z \textbackslash\{\}left( \textbackslash\{\}frac\{\textbackslash\{\}partial y\}\{\textbackslash\{\}partial z\}\textbackslash\{\}right)_x \textbackslash\{\}left( \textbackslash\{\}frac\{\textbackslash\{\}partial z\}\{\textbackslash\{\}partial x\}\textbackslash\{\}right)_y = -1. \textbackslash\{\}tag\{1\} \textbackslash\{\}]

Отметим также, что при дифференцировании при тех же зафиксированных переменных справедливо равенство

$$\textbackslash\{\}left( \textbackslash\{\}frac\{\textbackslash\{\}partial x\}\{\textbackslash\{\}partial y\}\textbackslash\{\}right)_z = \textbackslash\{\}frac\{1\}\{\textbackslash\{\}left( \{\textbackslash\{\}partial y\}/\{\textbackslash\{\}partial x\}\textbackslash\{\}right)_z\}.$$

\section*{Part A. Electrocaloric effect in ferroelectrics (5.5 points)}

In this part, the effect of changing the temperature of a ferroelectric under the influence of an electric field is studied from purely thermodynamic considerations.

\problembox{A1}{Consider one mole of an ideal gas. Considering its pressure $p$, volume $V$ and temperature $T$ as three quantities, check identity (1).}{0.50}

Further in this part we will consider a sample of a ferroelectric, which has the form of a thin plate placed between the plates of a flat capacitor. The volume of the ferroelectric $V$ can be considered constant, therefore its state is characterized by polarization $P$ and temperature $T$. A ferroelectric can be influenced by adding or removing heat or applying voltage (thereby creating an electric field).

\problembox{A2}{Let a cycle be performed with the ferroelectric (see Fig.), in coordinates $E$, $P$. Determine the total amount of heat received by the ferroelectric per cycle. Express your answer in terms of $E_1$, $E_2$, $P_1$, $P_2$, $V$.}{0.70}

Let a cycle be performed with the ferroelectric (see Fig.), in coordinates $E$, $P$. Determine the total amount of heat received by the ferroelectric per cycle. Express your answer in terms of $E_1$, $E_2$, $P_1$, $P_2$, $V$.

\begin{figure}[H]
    \centering
    \includegraphics[width=0.72\textwidth]{T4_EN_assets/fig2.jpg}
    \caption{Figure 2.}
\end{figure}

\problembox{A3}{Now consider the ABCD cycle (see figure), which has the form of an infinitesimal parallelogram.  Indicate the direction of traversing the cycle in which the amount of heat received by the ferroelectric during the cycle will be positive. Justify your answer.}{0.40}

Now consider the ABCD cycle (see figure), which has the form of an infinitesimal parallelogram.  Indicate the direction of traversing the cycle in which the amount of heat received by the ferroelectric during the cycle will be positive. Justify your answer.

\begin{figure}[H]
    \centering
    \includegraphics[width=0.72\textwidth]{T4_EN_assets/fig3.jpg}
    \caption{Figure 3.}
\end{figure}

The following parts of the problem are devoted to proving the identity \textbackslash\{\}[ \textbackslash\{\}left( \textbackslash\{\}frac\{\textbackslash\{\}partial S\}\{\textbackslash\{\}partial E\} \textbackslash\{\}right)_T = \textbackslash\{\}left( \textbackslash\{\}frac\{\textbackslash\{\}partial P\}\{\textbackslash\{\}partial T\}\textbackslash\{\}right)_E, \textbackslash\{\}tag\{2\} \textbackslash\{\}] where $S$ is the entropy of a unit volume of the ferroelectric.  To do this, consider the cycle from point $A3$, considering the sides of the parallelogram to be infinitesimal sections of isotherms and adiabats. Moreover, which of the sections are isotherms and which are adiabats depends on the parameters of the substance. Since the slope of section $BC$ is less than the slope $AB$, $BC$ will be an isotherm under the condition $$ \textbackslash\{\}left( \textbackslash\{\}frac\{\textbackslash\{\}partial P\}\{\textbackslash\{\}partial E\}\textbackslash\{\}right)_T < \textbackslash\{\}left( \textbackslash\{\}frac\{\textbackslash\{\}partial P\}\{\textbackslash\{\}partial E\}\textbackslash\{\}right)_S. $$ Будем расcmатривать различные случаи. Для определенности везде дальше считаем, что $(\partial P/\partial T)_E < 0$.

\problembox{A4}{Let $BC$ and $AD$ be isotherms, and in the section $BC$ the temperature of the ferroelectric is $T_1$, and in the section $AD$ – $T_2$, $|T_1 - T_2|\ll T_1$. Express the amount of heat $Q$ received by the ferroelectric for the entire cycle in terms of $V$, $(\partial P/\partial T)_E$ (derivative of polarization at a constant electric field), $T_1$, $T_2$ and the difference in electric fields at points $D$ and $A$: $d E = E_D- E_A \approx E_C - E_B$.}{0.50}

\problembox{A5}{Under the conditions of the previous paragraph, find the heat $Q_+$ that is supplied to the ferroelectric per cycle. Indicate which isotherm is where heat is supplied and which is removed.  Express your answer in terms of $V$, $T_1$, $T_2$, $\Delta E$, $(\partial S/\partial E)_T$ (derivative of entropy with respect to the electric field at constant temperature).}{0.50}

\problembox{A6}{Using the formula for the efficiency coefficient for the Carnot cycle, as well as the results $A4$, $A5$, prove relation (2). Consider that the total amount of heat received by the ferroelectric per cycle coincides with the useful work.}{0.30}

Let now the isotherms be the segments $AB$ (temperature $T_1$) and $CD$ (temperature $T_2$).

\problembox{A7}{Express the amount of heat $Q$ received by the ferroelectric for the entire cycle through the volume $V$, the change in the electric field on the isotherm $dE = E_B - E_A \approx E_C - E_D$, the temperatures $T_1$, $T_2$ and the partial derivatives $(\partial E/\partial T)_P$, $(\partial P/\partial E)_T$. Note: You may need to relate the change in electric field on the isotherm to the change in polarization using the partial derivative at constant temperature.}{0.60}

\problembox{A8}{Using identity (1) for the triple of variables $E$, $P$, $T$, express the heat $Q$ from the previous paragraph in terms of $V$, $T_1$, $T_2$, $dE$ and $(\partial P/\partial T)_E$.}{0.50}

\problembox{A9}{Write down an expression for the heat $Q_+$ supplied in the cycle. Express your answer in terms of $V$, $T_1$, $T_2$, $\Delta E$, $(\partial S/\partial E)_T$ (derivative of entropy with respect to the electric field at constant temperature).}{0.50}

\problembox{A10}{Using the results $A8$, $A9$ and the formula for the efficiency of the Carnot cycle, prove formula (2) for the case under consideration. The definition of useful work is the same as in $A6$.}{0.30}

\problembox{A11}{Using formula (2) (it can be applied without proof) and relation (1) for the triple of variables $S$, $T$, $E$, obtain the formula $$ \textbackslash\{\}left( \textbackslash\{\}frac\{\textbackslash\{\}partial T\}\{\textbackslash\{\}partial E\}\textbackslash\{\}right)_S  = - \textbackslash\{\}frac\{T\}\{C_E\} \textbackslash\{\}left(\textbackslash\{\}frac\{\textbackslash\{\}partial P\}\{\textbackslash\{\}partial T\} \textbackslash\{\}right)_E. $$Здесь $C_E$ – теплоемкость единицы объема сегнетоэлектрика при постоянном электрическом поле $$ C_E = \textbackslash\{\}left( \textbackslash\{\}frac\{\textbackslash\{\}delta Q\}\{dT\}\textbackslash\{\}right)_E. $$This formula shows that it is possible to change the temperature of the ferroelectric by adiabatically changing electric field.}{0.70}

\section*{Part B. Microscopic model of a ferroelectric (4.5 points)}

In order to achieve the strongest effect of changing the temperature of the ferroelectric, it is necessary to use a substance with a large value of the derivative $\left({\partial P}/{\partial T} \right)_E$. It turns out that this derivative is especially large near the phase transition point. As long as the ferroelectric temperature is above the phase transition temperature, its spontaneous polarization is zero, and when it is lower, spontaneous polarization occurs. Let us consider the following ferroelectric model. Let the ferroelectric contain ions of the crystal lattice, which are in the effective potential $$ U(x) = U_0 \textbackslash\{\}left( \textbackslash\{\}frac\{a_2\}\{2\} \textbackslash\{\}left( \textbackslash\{\}frac\{x\}\{x_0\}\textbackslash\{\}right)^2  - \textbackslash\{\}frac\{a_4\}\{4\} \textbackslash\{\}left( \textbackslash\{\}frac\{x\}\{x_0\}\textbackslash\{\}right)^4 + \textbackslash\{\}frac\{a_6\}\{6\} \textbackslash\{\}left( \textbackslash\{\}frac\{x\}\{x_0\}\textbackslash\{\}right)^6\textbackslash\{\}right) $$Для простоты расcmатривается одномерная задача, поэтому потенциал зависит только от одной координаты иона, которая отсчитывается от его положения равновесия. Здесь $U_0$ – характерное значение энергии, $x_0$ – характерное cmещение иона, $a_4 > 0$, $a_6 > 0$ – безразмерные постоянные, характеризующие форму потенциала, величина $a_2 = \alpha (T - T_с)/T_c$ ($\alpha$ – некоторый коэффициент) характеризует изменение потенциала с температурой. В зависимости от температуры у потенциала могут появляться дополнительные положения равновесия. При этом фактически реализуется то из устойчивых положений равновесия, которому отвечает наименьшее значение потенциала. Если ион находится в положении равновесия с $x = 0$, его вклад в электрический дипольный момент кристалла равен нулю. Если ион находится в положении равновесия с ненулевым $x$, у кристалла возникает спонтанная поляризация. Концентрация ионов в кристалле $n$, заряд одного иона $q$. Для удобства вычислений вы можете использовать безразмерную переменную $y = x/x_0$.

\problembox{B1}{Find the equilibrium positions of the ion in the potential. Express your answer in terms of $x_0$, $a_2$, $a_4$, $a_6$. Find the number of equilibrium positions depending on the value of $a_2$ and indicate which of them are stable.}{1.40}

\problembox{B2}{Determine numerically, starting from what value $a_{2c}$ the non-zero value $x$ corresponds to the minimum possible value of potential energy. Describe the method used. Use the numerical values ​​$a_4 = 1$, $a_6 = 1$.}{0.70}

\problembox{B3}{Find the dependence of the spontaneous polarization of the ferroelectric $P$ on temperature. Express the answer in terms of $n$, $q$, $x_0$, $a_4$, $a_6$, $\alpha$ and the dimensionless temperature $t = (T- T_c)/T_c$ Consider that as a result of interaction, all ions are shifted in the same direction during polarization.}{0.40}

\problembox{B4}{Using the results of the previous paragraphs, calculate the derivative of polarization with respect to temperature for an electric field equal to zero $(\partial P/\partial T)_{E = 0}$. Express your answer in terms of $n$, $q$, $x_0$, $a_4$, $a_6$, $\alpha$, $t$ and $T_c$. Plot the dependence of this derivative on $t$, indicate the characteristic points. When calculating point coordinates, use the numerical data from $B2$. The influence of the own electric field created by the ferroelectric on the ions can be neglected.}{0.70}

\problembox{B5}{Let now a weak external electric field $E$ be applied to the ferroelectric. Then the change in polarization of the dielectric can be considered linear in the field. Write down an equation from which you can find the equilibrium positions when the electric field is turned on. Find the polarizability from this equation $$ \textbackslash\{\}alpha = \textbackslash\{\}left(\textbackslash\{\}frac\{\textbackslash\{\}partial P\}\{\textbackslash\{\}partial E\} \textbackslash\{\}right)_\{T, E = 0\}. $$Выразите ответ через $n$, $q$, $x_0$, $U_0$, $a_4⟦M7 ⟧a_6$, $\alpha$, $t$ и положение равновесия иона без учета электрического поля $x$. Постройте график зависимости поляризуемости от обезразмеренной температуры $t$, укажите характерные точки. При расчете координат точек используйте численные данные из $B3$.}{1.30}

\vspace{1em}
\noindent\textit{Source:} \url{https://pho.rs/p/3851}

\end{document}